\section{Problem Description}
%\todo[inline,caption={}]{Describe what the premises and problems in the thesis are dealing with. 
%	Generally, this is often some lack of knowledge or lack of means that is planned to be resolved. For instance, non-existing insights in a certain domain or no tools/algorithms/implementations for further investigations.
%}

The proposed thesis is about formulating and evaluating a novel approach to online dynamic scene reconstruction via 3DGS that only looks at the subset of 3D Gaussians - the \textit{subscene} - that is being altered in comparison to the previous frame. Such a subscene contains only a fraction of all Gaussians and would therefore be much faster to compute.

In a first step, though, the space covered by the fixed multi-camera setup is trained as a regular 3DGS scene which functions as the base for subsequent updates, as with 4DGS.

The input imagery for a subscene are frame-to-frame pixel changes from the new frame to the previous frame (for each camera). This is explicitly not a difference image, as the input image would contain pixels that did change over two consecutive frames holding their changed RGB value. The rest of the pixels would be encoded as 'unset' and be ignored during subscene reconstruction.

During the forward pass of subscene reconstruction, only the set pixels (the ones not to be ignored) are considered and the loss function only compares the set pixels in rendered and input images before heading into the backward pass. Also in the backward pass, only Gaussians behind set pixels are considered - this is provided by 3DGS out of the box. All other Gaussians would not be touched during gradient descent (\enquote{frozen}).\vspace*{1em}

The algorithm is expected to converge very quickly as
\begin{itemize}
	\item the existing geometry is spatially close to what the new imagery of a new frame requires.
	%\item new geometry appearing where there was none close by in the scene before is most likely very small in the first frame it appears in, so it can just be reconstructed directly.
	\item solutions for new geometry entering the scene require spawning new 3D Gaussians at the correct position. This can be achieved by finding the intersection of changed pixels from at least two input images and mapping those to points in 3D space where the Gaussians can be spawned. We would spawn a larger amount of empty Gaussians in this spot and prune the ones that remained empty after each new frame. So there would be some kind of secondary pruning operation after the regular pruning during training of the base scene.
	\item geometry removed/leaving the scene will be handled in the regular way as the pixel color will then show the color behind the geometry that was removed. This geometry will have been known before and the algorithm should only remove the 3D Gaussians that were in front of the one that provided the color behind.
\end{itemize}
This means that the training algorithm for the subscene is to a large extent identical to regular 3DGS/4DGS, but using a mask on pixels and Gaussians to be considered.

This process could be sped up if Gaussians contained motion vectors making predicting the Gaussians' next positions more precise, reducing the amount of training iterations.

This subscene approach in a fixed multi-camera setup \st{has never been investigated but} looks very promising to at least reach a low number of FPS which can then later -- in subsequent work -- be refined to run at real-time speeds (\enquote{First make it work, then make it work fast}). Refer to Chapter \ref{ch:decogs} for an explanation of the strikethrough segment above.

\paragraph{Other Challenges}
\begin{itemize}
	\item Implementing such a niche idea cannot be done from scratch. As there are several open source repositories available, some time will have to be spent finding a suitable repository to base the implementation of the proposed thesis off of.
	\item In order to validate the implementation there are at least two datasets publicly available that have been used in previous papers for validation purposes. Alternatives would mean creating a scene from scratch through software like Blender.
	\item If there's time left, how could motion vectors be incorporated into the Gaussians?
	\item Maybe adjusting existing Gaussians and creating new Gaussians has to/should be done in two passes? How to classify existing vs. new geometry?
\end{itemize}



%For the real-time component we can deduce that we get 33 milliseconds per frame to advance the scene if we consider 30 frames per second as smooth enough. For a decently sized static initial scene that takes 5 minutes to train, each iteration of the training algorithm (there are usually around 30,000 iterations) takes 10ms. So the goal is to fill the available 33ms with efficient (GPU) logic to fully advance the scene by one frame. Seeing that the full scene takes 10ms for one training iteration the consequence is clear: what is being considered for the real-time calculation has to be much less than the full scene in order to work at all. However, in 33ms, a modern Nvidia consumer GPU like the GeForce 5090 can perform $3.45$ \textit{trillion} FP16 operations fully in parallel while datacenter GPUs like the Nvidia B200 can even reach 74.5 trillion operations (this is the theoretical maximum, 40-60\% of this are usually considered \enquote{good}).

%This is where the novel approach steps in. For the dynamic part of the 3D scene (the one after training the initial, static 3D scene) I propose to only use the differences between subsequent captured images of each camera to contribute to the scene. Fixed cameras will only ever see the same section of the scene, allowing only the changes (the \enquote{diff}) in the scene to propagate into the novel algorithm.

%We also store the previous diff and as a first step, those two diffs are compared and similar/same pixes are mapped onto each other.
%What the 3DGS algorithm does out of the box is map pixels to 3D Gaussians (1:n). So when we have a pixel, we know which 3D Gaussians will be affected. We mark all affected 3D Gaussians as \textit{active} and separate them from the full scene (+ some neighboring 3D Gaussians as padding
% and 3D Gaussians in the direction the pixel was moving
%). Then the centerpiece of the novel algorithm iterates over the diff-based input images to create a small, possibly tiny, new 3D scene, which will generate quickly as the active 3D Gaussians represent shapes and positions very close to the ones in the new frames. This small diff-based 3D scene is then inserted back into the full scene. As we only touched affected 3D Gaussians, this should not result in gaps etc.


%Geometry being removed/leaving the scene will be handled in the regular way as the pixel color will then show the color behind the geometry that was removed. This geometry will have been known before and the algorithm should only remove the 3D Gaussians that were in front of the one that provided the color behind.

%learning directly from encoded sparse image?

%RLE?




%20 minutes = 1200000 ms, divided by 30000 iterations = 40ms per iteration

%5 minutes = 10ms per iteration

%Look into SLAM applications

%Für weitere Informationen zum Darstellen von Argumentationsgraphen, siehe\footnote{\tiny\ttfamily\url{https://github.com/aig-hagen/tikz_argumentation/blob/main/argumentation-doc.pdf}}.

% \begin{table*}[t]
% \begin{center}
% \begin{tabular}{lll}
% \hline
% 1&2&3\\
% 4&5&6\\
% \hline
% \end{tabular}
% \end{center}
% \caption{Eine Tabelle.} \label{tab:t1}
% \end{table*}


% For propositions and definitions please use appropriate environments.
% Here is an example for a definition.

% \begin{definition}
% 	A deterministic finite automaton (DFA) is a tuple \( A=(Q,\Sigma,\delta,q_0,F) \) consisting of
% 	\begin{itemize}
% 		\item a finite set of states \( Q \),
% 		\item a finite alphabet \( \Sigma \),
% 		\item a transition function \( \delta:Q\times\Sigma\to Q \), 
% 		\item a start state \( q_0 \in Q \), and
% 		\item a set of accepting states \( F \subseteq Q \).
% 	\end{itemize}
% \end{definition}

% Here is an example for a proposition.

% \begin{proposition}
% 	The problem of deciding whether a Turing machine halts for an input string is undecidable.
% \end{proposition}

% Examples can be made with the example environment.

% \begin{example}
% 	Content of the example environment.
% \end{example}